\section{Fundamentos teóricos}

La resistencia a compresión del concreto es una respuesta cuantitativa que
depende de la composición de la mezcla y de su edad. El conjunto Concrete
Compressive Strength reúne ocho predictoras: cemento, escoria de alto horno,
ceniza volante, agua, superplastificante, agregado grueso, agregado fino y edad.
Los componentes se expresan en kg/m$^3$, la edad en días y la resistencia en MPa.
Esta estructura define un problema de regresión supervisada: aprender una
función que estime una cantidad a partir de variables observadas. La relación
entre ingredientes y resistencia puede ser no lineal, lo que fundamenta la
comparación de modelos con distintas capacidades de representación
\cite{uci_concrete,yeh1998}.

La regresión lineal representa la respuesta mediante
$\hat y=\beta_0+\sum_{j=1}^{p}\beta_jx_j$. El modelo simple utiliza una
predictora y el múltiple incorpora varias. Los coeficientes de mínimos
cuadrados minimizan la suma de residuos al cuadrado y proporcionan una
referencia interpretable. Cuando las predictoras están relacionadas, las
estimaciones pueden ser sensibles a variaciones de los datos. Ridge introduce
una penalización $\ell_2$ que contrae los coeficientes; Lasso utiliza una
penalización $\ell_1$ capaz de producir coeficientes exactamente nulos
\cite{sklearn_lineales,hoerl1970,tibshirani1996}.

El parámetro $\alpha$ controla la intensidad de regularización. LassoCV y
RidgeCV seleccionan este parámetro mediante validación cruzada y pertenecen a
las mismas familias que Lasso y Ridge, respectivamente. Comparar estas
variantes permite examinar el efecto de la regularización y de su selección
automática. Los nueve identificadores evaluados representan algoritmos o
variantes; no constituyen nueve familias conceptuales independientes. El
benchmarking distingue la comparación dentro de una familia de la comparación
entre representaciones lineales, vectores de soporte y árboles
\cite{sklearn_lineales,cawley2010}.

La regresión por vectores de soporte combina control de complejidad con una
pérdida que tolera desviaciones dentro de un margen $\epsilon$. El parámetro
$C$ regula el compromiso entre penalizar errores y controlar la complejidad
del ajuste. Los kernels permiten representar relaciones que no son lineales
en las variables originales. Por ello, el desempeño de SVR depende del kernel,
de los hiperparámetros y de la escala de las predictoras. Su comparación
requiere las mismas particiones y un protocolo de evaluación común
\cite{smola2004}.

Un árbol de regresión divide el espacio de entrada mediante reglas y produce
una aproximación constante por regiones. Su profundidad controla parte de su
flexibilidad: un árbol muy complejo puede ajustar particularidades del
entrenamiento y generalizar peor. Random Forest agrega predicciones de
árboles aleatorizados; el promedio busca reducir la variabilidad del estimador
individual. Comparar ambos permite estudiar el efecto del ensamble y
contrastar estas representaciones con los modelos lineales y basados en kernels
\cite{sklearn_arboles,breiman2001}.

La evaluación utiliza MAE, MSE, RMSE y $R^2$. Para $n$ observaciones con
respuestas $y_i$, predicciones $\hat y_i$ y media observada $\bar y$, se definen
\cite{sklearn_metricas}:
\begin{align*}
\mathrm{MAE} &= \frac{1}{n}\sum_{i=1}^{n}|y_i-\hat y_i|,\\
\mathrm{MSE} &= \frac{1}{n}\sum_{i=1}^{n}(y_i-\hat y_i)^2,\\
\mathrm{RMSE} &= \sqrt{\frac{1}{n}\sum_{i=1}^{n}(y_i-\hat y_i)^2},\\
R^2 &= 1-\frac{\sum_{i=1}^{n}(y_i-\hat y_i)^2}
{\sum_{i=1}^{n}(y_i-\bar y)^2}.
\end{align*}
El MAE resume la magnitud absoluta del error; el MSE aumenta la influencia de
desviaciones grandes al elevarlas al cuadrado. MAE y RMSE conservan la unidad
de la respuesta, MPa, mientras MSE se expresa en MPa$^2$. $R^2$ puede ser
negativo y no representa un porcentaje de aciertos; su expresión requiere
variación en la respuesta del conjunto evaluado. La selección utiliza RMSE
de validación y la prueba se reserva para evaluar los modelos seleccionados.
El preprocesamiento aprendido debe ajustarse con entrenamiento para evitar
que información de evaluación influya en el modelo. Incluso con esta
separación, diferencias pequeñas en una muestra no acreditan superioridad
general entre algoritmos \cite{sklearn_fuga,cawley2010}.

\begin{figure}[H]
\centering
\begin{tikzpicture}[
  caja/.style={draw=black!55,rounded corners=2pt,align=center,
    inner sep=6pt,font=\small},
  flecha/.style={-{Stealth[length=2mm]},thick,draw=black!65}
]
\node[caja,fill=gray!10,text width=6.2cm] (datos) at (0,0)
  {\textbf{Regresión supervisada}\\Predictoras cuantitativas
  $\longrightarrow$ resistencia a compresión};
\node[caja,fill=gray!10,text width=10.7cm,below=0.55cm of datos] (preparacion)
  {\textbf{Protocolo común}\\Partición de datos y preprocesamiento
  ajustado con entrenamiento};
\node[caja,fill=blue!8,text width=3.1cm] (lineales) at (-4.2,-3.3)
  {\textbf{Modelos lineales}\\Simple y múltiple\\Lasso y LassoCV\\Ridge y RidgeCV};
\node[caja,fill=green!9,text width=3.1cm] (soporte) at (0,-3.3)
  {\textbf{Vectores de soporte}\\SVR\\Kernel, $C$ y $\epsilon$};
\node[caja,fill=orange!12,text width=3.1cm] (arboles) at (4.2,-3.3)
  {\textbf{Árboles y ensambles}\\Árbol de decisión\\Random Forest};
\node[caja,fill=gray!10,text width=10.7cm] (validacion) at (0,-5.3)
  {\textbf{Comparación por RMSE de validación}\\Dentro de familias y
  entre familias; selección de configuraciones};
\node[caja,fill=gray!10,text width=7.1cm,below=0.55cm of validacion] (prueba)
  {\textbf{Evaluación en prueba reservada}\\MAE, MSE, RMSE y $R^2$;
  residuos, interpretación y límites};
\draw[flecha] (datos) -- (preparacion);
\draw[flecha] (preparacion.south) -- (lineales.north);
\draw[flecha] (preparacion.south) -- (soporte.north);
\draw[flecha] (preparacion.south) -- (arboles.north);
\draw[flecha] (lineales.south) -- (validacion.north);
\draw[flecha] (soporte.south) -- (validacion.north);
\draw[flecha] (arboles.south) -- (validacion.north);
\draw[flecha] (validacion) -- (prueba);
\end{tikzpicture}
\caption{Mapa conceptual del benchmarking de regresión: familias,
preparación, selección por validación y evaluación en prueba.}
\label{fig:mapa-regresion}
\end{figure}
